% Vietnamese translation for OI Public Library
% Source-File: IOI中国国家候选队论文/2006/王栋/王栋.ppt
\documentclass[11pt]{article}
\usepackage{oipl-vn}
\usepackage{tikz}
\usetikzlibrary{arrows.meta,calc}

\newcommand{\Oh}{\mathcal{O}}
\newcommand{\Vor}{\operatorname{Vor}}
\newcommand{\DT}{\operatorname{DT}}

\title{Voronoi\\{\large Ghi chú theo slide}}
\author{Wang Dong}
\date{2006}

\begin{document}
\maketitle

\origtitle{Voronoi diagrams}
\sourcepath{IOI China candidate team papers / 2006 / Wang Dong / Wang Dong.ppt}

\section{Phạm vi}

Bộ slide đi kèm nguồn về biểu đồ Voronoi. Tệp DOC cùng chủ đề bị hỏng một
phần, nên bản ghi này dựa trên bản dịch phục dựng từ các mảnh văn bản, slide
và mã Pascal đi kèm.

\section{Ô Voronoi}

Với tập điểm \(S\), ô Voronoi của điểm \(P_i\) gồm các điểm trên mặt phẳng gần
\(P_i\) hơn mọi điểm còn lại. Biên giữa hai ô nằm trên đường trung trực của
hai điểm sinh ô. Tập các ô tạo thành biểu đồ \(\Vor(S)\).

\begin{figure}[ht]
\centering
\begin{tikzpicture}[scale=0.95, site/.style={circle, fill=black, inner sep=1.8pt}]
  \coordinate (a) at (0,0);
  \coordinate (b) at (4,0);
  \coordinate (c) at (1.5,2.4);
  \coordinate (d) at (3.3,2.0);
  \draw[gray!55, dashed] (a) -- (b) -- (d) -- (c) -- (a) -- (d);
  \draw[blue!70, thick] (2,-1.0) -- (2,2.8);
  \draw[blue!70, thick] (-0.2,1.15) -- (3.2,-0.95);
  \draw[blue!70, thick] (0.9,1.9) -- (4.5,0.7);
  \node[site, label=below:\(P_i\)] at (a) {};
  \node[site, label=below:\(P_j\)] at (b) {};
  \node[site, label=above:\(P_k\)] at (c) {};
  \node[site, label=above:\(P_\ell\)] at (d) {};
  \draw[-Latex] (2.45,1.5) -- (2.05,1.5)
    node[midway, above] {biên Voronoi};
  \node[gray!70] at (2,-0.35) {trung trực của \(P_iP_j\)};
  \node[gray!70, align=center] at (3.35,2.65)
    {cạnh Delaunay\\nối hai ô kề};
\end{tikzpicture}
\caption{Biên Voronoi là trung trực giữa hai điểm sinh ô; cạnh Delaunay nối các điểm có ô kề nhau.}
\end{figure}

\section{Đối ngẫu Delaunay}

Tam giác hóa Delaunay \(\DT(S)\) là đồ thị đối ngẫu của biểu đồ Voronoi: hai
điểm được nối nếu hai ô Voronoi tương ứng kề nhau. Một tam giác Delaunay có
tính chất đường tròn ngoại tiếp không chứa điểm nào khác của \(S\) ở bên
trong. Quan hệ đối ngẫu này thường giúp chuyển bài toán khoảng cách cực trị
sang xét cạnh hoặc đỉnh hữu hạn.

\section{Dựng và ứng dụng}

Slide và mã đi kèm gợi ý hai hướng dựng: giao các nửa mặt phẳng xác định ô
Voronoi, hoặc dựng thông qua tam giác hóa Delaunay. Các ví dụ \emph{Run Away}
và \emph{Fat Man} đều dùng ý tưởng rằng điểm tối ưu theo khoảng cách tới các
vật cản thường nằm ở đỉnh Voronoi, trên cạnh biên, hoặc tại điểm đặc biệt của
miền xét.

\section{Cài đặt}

Cần xử lý cẩn thận giao đường thẳng, so sánh phía của điểm, trường hợp nhiều
điểm đồng viên và sai số số thực. Trong bài thi, phần hình học không chỉ là
công thức; thứ tự xét biên và quy ước sai số quyết định độ ổn định của chương
trình.

\end{document}
